\paragraph{A claim needs evidence as well as a place in memory.}
The preceding lifecycle explains what an identified sale does to the book.
It does not establish where the invoice identity, amount or due date came
from. An implementer needs to preserve the source reading separately from the
decision to recognize a sale. Otherwise a correction to a document can be
mistaken for a second transaction, or a missing payment term can acquire an
invented value merely because execution requires one.

The bounded \texttt{ReceivableSourceContract/1} supplies this separation for
explicitly constructed JSON documents. Each reading binds an exact UTF-8 byte
span and its hash, the company/book/currency context, an assertion identity,
and a stable claim and sale-event identity. Original and corrected assertions
remain present. A named information cutoff controls which assertion was
available; a separate ingestion cutoff controls which version the system had
received. Unknown timing, competing corrections, altered source bytes or
unreadable terms remain unresolved or refuse as appropriate. Recognition
approval, receipts and assessment policy are separately supplied premises.

Extend the synthetic company by a second sale on 3~January 2027: consideration
of USD~40, cost of USD~10, due on 28~February, to the same named customer but
under a different invoice identity. After the earlier USD~10 receipt, the two
claims are USD~20 and USD~40. Cash is USD~110, inventory is zero, profit is
USD~50 and assets and equity are USD~170. Carry preserves both claims. The
subsequent receipts of USD~20 and USD~40 raise cash to USD~170 and clear the
claims, with no new profit and operating cash flow of USD~60. These balances
follow from the declared transactions; they are not readings of a real company.

Now suppose a corrected assertion changes the first invoice's due date to
31~March. It becomes available at noon UTC on 15~March and is ingested at
09:00 UTC on 16~March. An earlier information or ingestion view retains the
original assertion. A later view selects the correction and reconstructs a
separate history from the same opening. Its February overdue assessment is
zero instead of USD~20; its journals and financial statements are unchanged.
Three assertions therefore still describe two claims. This operation corrects
a reading of the original terms. A genuine renegotiation with its own effective
date would require an amendment method; it cannot be represented by silently
rewriting the original sale.

\paragraph{A disclosed receivable may have an unidentified population.}
The retained Apple statements expose the next limit. Accounts receivable, net
are displayed as USD~33,410 million at 28~September 2024 and USD~29,639 million
at 28~December 2024. These are consolidated measurements with displayed
precision, not gross invoices with customer identities and payment dates. The
annual cash-flow entry of USD~(3,788) million is a change in operating assets;
it is not an identified customer collection. The receivable notes provide
concentration information, but that information does not supply a complete
invoice schedule.\footnote{Retained annual Form 10-K, PDF pages 34, 36 and 41
(printed pages 31, 33 and 38); retained interim Form 10-Q, balance sheet and
Note 4, Accounts Receivable. Original observations, reading-file entries and
source hashes are retained in
\path{data/accounting-master/retained-receivable-boundary-20261025/}.}
\citep{apple2024,apple2025q1}

The distinction also affects method choice. Apple's vendor non-trade
receivables arise from component sales to manufacturing vendors. The disclosed
treatment excludes those sales from products net sales and recognizes gains
as reductions in products cost of sales when the related final products are
sold. The generic credit-sale lifecycle above recognizes revenue at its sale
event, so applying it to these vendor transactions would compute a different
treatment. Preserving the vendor caption is useful even while the required
transaction facts and accounting method remain unspecified.

For such observations, \texttt{SourceClaimPopulationView/1} attaches the
existing \texttt{MethodCollections/1} partial collection to the unchanged
source-constrained company state. A partial collection can have no observed
rows while the actual number of claims remains unknown. Its claim count,
gross remainder and contractual due terms return explicit unknown results.
This differs from an explicitly complete empty collection, whose count and
total are zero. A complete constructed collection containing two uniquely
identified remainders of USD~20 and USD~40 has count two and total USD~60,
conditional on the supplied completeness premise. The retained-source adapter
does not accept an override that could manufacture that premise.

The reported measurement remains available in all these cases, together with
its source context, date, precision and unresolved obligations. An explicitly
assumed rounding interval can constrain that same measurement; it supplies
neither members nor a gross-to-net reconciliation. No balancing residual or
legal-book allocation is inferred. Even a known contractual due date does
not establish when a customer will pay: settlement behavior still needs its
own evidence or scenario assumption. The attachment consequently exposes no
executable opening and authorizes no posting.

The source-contract, retained-disclosure and partial-population references
preserve complete typed requests, results and replay checks. Their reference
directories are listed below under the common
\path{data/accounting-master/} parent.
\begin{flushleft}\small\ttfamily
populated-source-contracts-20261025/\\
retained-receivable-boundary-20261025/\\
partial-claim-populations-20261025/
\end{flushleft}
These references connect readings to useful conditional calculations without treating missing claim
state as an empty book. Real-source admission, complete opening populations,
impairment, transaction-specific recognition and broader company composition
remain substantive requirements.
